\Chapter{50}

\Verse{1}
{
    \zA{话说薛宝钗道：“到底分个次序，让我写出来。”}
}
{
    \eA{But to continue. "We should, after all," Pao-ch'ai suggested, "make some
        distinction as to order.}
}
{
}

\Verse{2}
{
    \zA{说着，便令众人拈阄为序。}
}
{
    \eA{Let me write out what's needful." After uttering this proposal,}
}
{
}

\Verse{3}
{
    \zA{起首恰是李氏，然后按次各各开出。}
}
{
    \eA{she urged every one to draw lots and determine the precedence. The first one
        to draw was Li Wan.}
}
{
}

\Verse{4}
{
    \zA{凤姐儿道：“既这么说，我也说一句在上头。”}
}
{
    \eA{After her, a list of the respective names was made in the order in which
        they came out. "Well, in that case," lady Feng rejoined,}
}
{
}

\Verse{5}
{
    \zA{众人都笑起来了，说：“这么更妙了。”}
}
{
    \eA{"I'll also give a top line." The whole party laughed in chorus. "It will be
        ever so much better like this,"}
}
{
}

\Verse{6}
{
    \zA{宝钗将“稻香老农”之上补了一个“凤” 字，李纨又将题目讲给他听。}
}
{
    \eA{they said. Pao-ch'ai supplied above 'the old labourer of Tao Hsiang' the
        word 'Feng,' whereupon Li Wan went on to explain the theme to her. "You
        musn't poke fun at me!"}
}
{
}

\Verse{7}
{
    \zA{凤姐儿想了半天，笑道：“你们别笑话我，我只有了 一句粗话，可是五个字的。}
}
{
    \eA{lady Feng smiled, after considerable reflection. "I've only managed to get a
        coarse line. It consists of five words. As for the rest, I have no idea how
        to manage them."}
}
{
}

\Verse{8}
{
    \zA{下剩的我就不知道了。”}
}
{
    \eA{"The coarser the language, the better it is," one and all laughed.}
}
{
}

\Verse{9}
{
    \zA{众人都笑道：“越是粗话越好。}
}
{
    \eA{"Out with it! You can then go and attend to your legitimate business!"}
}
{
}

\Verse{10}
{
    \zA{你说了，就只管干正事去罢。”}
}
{
    \eA{"I fancy," lady Feng observed, "that when it snows there's bound to be
        northerly wind,}
}
{
}

\Verse{11}
{
    \zA{凤姐儿笑道：“想下雪必刮北风，昨夜听见一夜的 北风，我有一句，这一句就是‘一夜北风紧’。}
}
{
    \eA{for last night I heard the wind blow from the north the whole night long.
        I've got a line, it's: "'The whole night long the northern wind was high;'
        "but whether it will do or not, I am not going to worry my mind about it."
        One and all, upon hearing this, exchanged looks. "This line is,}
}
{
}

\Verse{12}
{
    \zA{使得使不得，我就不管了。”}
}
{
    \eA{it's true, coarse," they smiled, "and gives no insight into what comes
        below,}
}
{
}

\Verse{13}
{
    \zA{众人 听说，都相视笑道：“这句虽粗，不见底下的，这正是会作诗的起法。}
}
{
    \eA{but it's just the kind of opening that would be used by such as understand
        versification. It's not only good, but it will afford to those, who come
        after you, inexhaustible scope for writing. In fact, this line will take the
        lead,}
}
{
}

\Verse{14}
{
    \zA{不但好，而 且留了写不尽的多少地步与后人。}
}
{
    \eA{so 'old labourer of Tao Hsiang' be quick and indite some more to tag on
        below." Lady Feng, 'sister-in-law' Li, and P'ing Erh had then another couple
        of glasses,}
}
{
}

\Verse{15}
{
    \zA{就是这句为首，稻香老农快写上，续下去。”}
}
{
    \eA{after which each went her own way. During this while Li Wan wrote down: The
        whole night long the northern wind was high;}
}
{
}

\Verse{16}
{
    \zA{凤 姐儿和李婶娘平儿又吃了两杯酒，自去了。}
}
{
    \eA{and then she herself subjoined the antithetical couplet: The door I ope, and
        lo the flakes of snow are still toss'd by the wind,}
}
{
}

\Verse{17}
{
    \zA{这里李纨就写了： 一夜北风紧， 自己联道： 开门雪尚飘。}
}
{
    \eA{And drop into the slush. Oh, what a pity they're so purely white! Hsiang
        Ling recited: All o'er the ground is spread, alas, this bright, refulgent
        gem;}
}
{
}

\Verse{18}
{
    \zA{入泥怜洁白， 香菱道： 匝地惜琼瑶。}
}
{
    \eA{But with an aim; for it is meant dry herbage to revive. T'an Ch'un said:
        Without design the dying sprouts of grain it nutrifies.}
}
{
}

\Verse{19}
{
    \zA{有意荣枯草， 探春道： 无心饰萎苗。}
}
{
    \eA{But in the villages the price of mellow wine doth rise. Li Ch'i added: In a
        good year, grain in the house is plentiful.}
}
{
}

\Verse{20}
{
    \zA{价高村酿熟， 李绮道： 年稔府粱饶。}
}
{
    \eA{The bulrush moves and the ash issues from the tube. Li Wen continued: What
        time spring comes the handle of the Dipper turns.}
}
{
}

\Verse{21}
{
    \zA{葭动灰飞管， 李纹道： 阳回斗转杓。}
}
{
    \eA{The bleaky hills have long ago their verdure lost. Chou-yen proceeded: On a
        frost-covered stream,}
}
{
}

\Verse{22}
{
    \zA{寒山已失翠， 岫烟道： 冻浦不生潮。}
}
{
    \eA{no tide can ever rise. Easy the snow hangs on the sparse-leaved willow
        twigs.}
}
{
}

\Verse{23}
{
    \zA{易挂疏枝柳， 湘云道： 难堆破叶蕉。}
}
{
    \eA{Hsiang-yün pursued: Hard 'tis for snow to pile on broken plantain leaves.
        The coal, musk-scented, burns in the precious tripod.}
}
{
}

\Verse{24}
{
    \zA{麝煤融宝鼎， 宝琴道： 绮袖笼金貂。}
}
{
    \eA{Pao-ch'in recited: Th' embroidered sleeve enwraps the golden sable in its
        folds.}
}
{
}

\Verse{25}
{
    \zA{光夺窗前镜， 黛玉道： 香粘壁上椒。}
}
{
    \eA{The snow transcends the mirror by the window in lustre. Pao-yü suggested:
        The fragrant pepper clings unto the wall.}
}
{
}

\Verse{26}
{
    \zA{斜风仍故故， 宝玉道： 清梦转聊聊。}
}
{
    \eA{The side wind still in whistling gusts doth blow. Tai-yü added: A quiet
        dream becomes a cheerless thing.}
}
{
}

\Verse{27}
{
    \zA{何处梅花笛？}
}
{
    \eA{Where is the fife with plum bloom painted on?}
}
{
}

\Verse{28}
{
    \zA{宝钗道： 谁家碧玉箫？}
}
{
    \eA{Pao-ch'ai continued: In whose household is there a flute made of green jade?}
}
{
}

\Verse{29}
{
    \zA{鳌愁坤轴陷， 李纨笑道：“我替你们看热酒去罢。”}
}
{
    \eA{The fish fears lest the earth from its axis might drop. "I'll go and see
        that the wine is warm for you people,"}
}
{
}

\Verse{30}
{
    \zA{宝钗命宝琴续联，只见湘云起来道： 龙斗阵云销。}
}
{
    \eA{Li Wan smiled. But when Pao-ch'ai told Pao-ch'in to connect some lines, she
        caught sight of Hsiang-yün rise to her feet and put in: What time the dragon
        wages war,}
}
{
}

\Verse{31}
{
    \zA{野岸回孤棹， 宝琴也联道： 吟鞭指灞桥。}
}
{
    \eA{the clouds dispel. Back to the wild shore turns the man with single scull.
        Pao-ch'in thereupon again appended the couplet:}
}
{
}

\Verse{32}
{
    \zA{赐裘怜抚戍， 湘云那里肯让人？}
}
{
    \eA{The old man hums his lines, and with his whip he points at the 'Pa' bridge.
        Fur coats are, out of pity,}
}
{
}

\Verse{33}
{
    \zA{且别人也不如他敏捷，都看他扬眉挺身的说道： 加絮念征徭。}
}
{
    \eA{on the troops at the frontiers bestowed. But would Hsiang-yün allow any one
        to have a say? The others could not besides come up to her in quickness of
        wits so that, while their eyes were fixed on her,}
}
{
}

\Verse{34}
{
    \zA{坳垤审夷险， 宝钗连声赞好，也便联道： 枝柯怕动摇。}
}
{
    \eA{she with eyebrows uplifted and figure outstretched proceeded to say: More
        cotton coats confer, for bear in memory th' imperial serfs!}
}
{
}

\Verse{35}
{
    \zA{皑皑轻趁步， 黛玉忙联道： 剪剪舞随腰。}
}
{
    \eA{The rugged barbarous lands are (on account of snow) with dangers fraught.
        Pao-ch'ai praised the verses again and again,}
}
{
}

\Verse{36}
{
    \zA{苦茗成新赏， 一面说，一面推宝玉命他联。}
}
{
    \eA{and next contributed the distich: The twigs and branches live in fear of
        being tossed about. With what whiteness and feath'ry step the flakes of snow
        descend!}
}
{
}

\Verse{37}
{
    \zA{宝玉正看宝琴、宝钗、黛玉三人共战湘云，十分有趣， 那里还顾得联诗？}
}
{
    \eA{Tai-yü eagerly subjoined the lines: The snow as nimbly falls as moves the
        waist of the 'Sui' man when brandishing the sword. The tender leaves of tea,
        so acrid to the taste, have just been newly brewed and tried.}
}
{
}

\Verse{38}
{
    \zA{今见黛玉推他，方联道： 孤松订久要。}
}
{
    \eA{As she recited this couplet, she gave Pao-yü a shove and urged him to go on.
        Pao-yü was, at the moment,}
}
{
}

\Verse{39}
{
    \zA{泥鸿从印迹， 宝琴接着联道： 林斧或闻樵。}
}
{
    \eA{enjoying the intense pleasure of watching the three girls Pao-ch'ai,
        Pao-ch'in and Tai-yü make a joint onslaught on Hsiang-yün,}
}
{
}

\Verse{40}
{
    \zA{伏象千峰凸， 湘云忙联道： 盘蛇一径遥。}
}
{
    \eA{so that he had of course not given his mind to tagging any antithetical
        verses. But when he now felt Tai-yü push him he at length chimed in with:}
}
{
}

\Verse{41}
{
    \zA{花缘经冷结， 宝钗和众人又都赞好，探春联道： 色岂畏霜凋。}
}
{
    \eA{The fir is the sole tree which is decreed for ever to subsist. The wild
        goose follows in the mud the prints and traces of its steps. Pao-ch'in took
        up the clue,}
}
{
}

\Verse{42}
{
    \zA{深院惊寒雀， 湘云正渴了，忙忙的吃茶，已被岫烟抢着联道： 空山泣老？。}
}
{
    \eA{adding: In the forest, the axe of the woodcutter may betimes be heard. With
        (snow) covered contours, a thousand peaks their heads jut in the air.}
}
{
}

\Verse{43}
{
    \zA{阶墀随上下， 湘云忙丢了茶杯联道： 池水任浮漂。}
}
{
    \eA{Hsiang-yün with alacrity annexed the verses: The whole way tortuous winds
        like a coiled snake. The flowers have felt the cold and ceased to bud.}
}
{
}

\Verse{44}
{
    \zA{照耀临清晓， 黛玉忙联道： 缤纷入永宵。}
}
{
    \eA{Pao-ch'ai and her companions again with one voice eulogised their fine
        diction. T'an Ch'un then continued:}
}
{
}

\Verse{45}
{
    \zA{诚忘三尺冷， 湘云忙笑联道： 瑞释九重焦。}
}
{
    \eA{Could e'er the beauteous snow dread the nipping of frost? In the deep court
        the shivering birds are startled by its fall.}
}
{
}

\Verse{46}
{
    \zA{僵卧谁相问， 宝琴也忙笑联道： 狂游客喜招。}
}
{
    \eA{Hsiang-yün happened to be feeling thirsty and was hurriedly swallowing a cup
        of tea, when her turn was at once snatched by Chou-yen,}
}
{
}

\Verse{47}
{
    \zA{天机断缟带， 湘云又忙道： 海市失鲛绡。}
}
{
    \eA{who gave out the lines, On the bare mountain wails the old man Hsiao. The
        snow covers the steps,}
}
{
}

\Verse{48}
{
    \zA{黛玉不容他道出，接着便道： 寂寞封台榭， 湘云忙联道： 清贫怀箪瓢。}
}
{
    \eA{both high and low. Hsiang-yün immediately put away the tea-cup and added: On
        the pond's surface, it allows itself to float. At the first blush of dawn
        with effulgence it shines.}
}
{
}

\Verse{49}
{
    \zA{宝琴也不容情，也忙道： 烹茶水渐沸， 湘云见这般，自为得趣，又是笑，又忙联道： 煮酒叶难烧。}
}
{
    \eA{Tai-yü recited with alacrity the couplet: In confused flakes, it ceaseless
        falls the whole night long. Troth one forgets that it implies three feet of
        cold. Hsiang-yün hastened to smilingly interpose with the distich:}
}
{
}

\Verse{50}
{
    \zA{黛玉也笑道： 没帚山僧扫， 宝琴也笑道： 埋琴稚子挑。}
}
{
    \eA{Its auspicious descent dispels the Emperor's grief. There lies one
        frozen-stiff, but who asks him a word? Pao-ch'in too speedily put on a smile
        and added:}
}
{
}

\Verse{51}
{
    \zA{湘云笑弯了腰，忙念了一句，众人问道：“到底说的是什么？”}
}
{
    \eA{Glad is the proud wayfarer when he's pressed to drink. Snapped is the
        weaving belt in the heavenly machine. Hsiang-yün once again eagerly quoted
        the line:}
}
{
}

\Verse{52}
{
    \zA{湘云道： 石楼闲睡鹤， 黛玉笑得握着胸口，高声嚷道： 锦？}
}
{
    \eA{In the seaside market is lost a silk kerchief. But Lin Tai-yü would not let
        her continue, and taking up the thread, she forthwith said: With quiet
        silence,}
}
{
}

\Verse{53}
{
    \zA{暖亲猫。}
}
{
    \eA{it enshrouds the raiséd kiosque.}
}
{
}

\Verse{54}
{
    \zA{宝琴也忙笑道： 月窟翻银浪， 湘云忙联道： 霞城隐赤标。}
}
{
    \eA{Hsiang-yün vehemently gave the antithetical verse: The utter poor clings to
        his pannier and his bowl. Pao-ch'in too would not give in as a favour to any
        one,}
}
{
}

\Verse{55}
{
    \zA{黛玉忙笑道： 沁梅香可嚼， 宝钗笑称：“好句！”}
}
{
    \eA{so hastily she exclaimed: The water meant to brew the tea with gently
        bubbles up. Hsiang-yün saw how excited they were getting and she thought it
        naturally great fun.}
}
{
}

\Verse{56}
{
    \zA{也忙联道： 淋竹醉堪调。}
}
{
    \eA{Laughing, she eagerly gave out: When wine is boiled with leaves 'tis not
        easy to burn.}
}
{
}

\Verse{57}
{
    \zA{宝琴也忙道： 或湿鸳鸯带， 湘云忙联道： 时凝翡翠翘。}
}
{
    \eA{Tai-yü also smiled while suggesting: The broom, with which the bonze
        sweepeth the hill, is sunk in snow. Pao-ch'in too smilingly cried:}
}
{
}

\Verse{58}
{
    \zA{黛玉又忙道： 无风仍脉脉， 宝琴又忙笑联道： 不雨亦潇潇。}
}
{
    \eA{The young lad takes away the lute interred in snow. Hsiang-yün laughed to
        such a degree that she was bent in two; and she muttered a line with such
        rapidity that one and all inquired of her:}
}
{
}

\Verse{59}
{
    \zA{湘云伏着，已笑软了。}
}
{
    \eA{"What are you, after all, saying?" In the stone tower leisurely sleeps the
        stork.}
}
{
}

\Verse{60}
{
    \zA{众人看他三人对抢，也都不顾作诗，看着也只是笑。}
}
{
    \eA{Hsiang-yün repeated. Tai-yü clasped her breast so convulsed was she with
        laughter. With loud voice she bawled out: Th' embroidered carpet warms the
        affectionate cat.}
}
{
}

\Verse{61}
{
    \zA{黛玉还 推他往下联，又道：“你也有才尽力穷之时!我听听，还有什么舌头嚼了？”}
}
{
    \eA{Pao-ch'in quickly, again laughingly, exclaimed: Inside Selene's cave lo,
        roll the silvery waves. Hsiang-yün added, with eager haste: Within the city
        walls at eve was hid a purple flag. Tai-yü with alacrity continued with a
        smile:}
}
{
}

\Verse{62}
{
    \zA{湘云 只伏在宝钗怀里笑个不住。}
}
{
    \eA{The fragrance sweet, which penetrates into the plums, is good to eat.
        Pao-ch'ai smiled.}
}
{
}

\Verse{63}
{
    \zA{宝钗推他起来，道：“你有本事，把‘二萧’的韵全用 完了，我才服你。”}
}
{
    \eA{"What a fine line!" she ejaculated; after which, she hastened to complete
        the couplet by saying: The drops from the bamboo are meet, when one is
        drunk, to mix with wine.}
}
{
}

\Verse{64}
{
    \zA{湘云起身笑道：“我也不是作诗，竟是抢命呢！”}
}
{
    \eA{Pao-ch'in likewise made haste to add: Betimes, the hymeneal girdle it
        moistens. Hsiang-yün eagerly paired it with:}
}
{
}

\Verse{65}
{
    \zA{众人笑道： “倒是你自己说罢。”}
}
{
    \eA{Oft, it freezeth on the kingfisher shoes. Tai-yü once more exclaimed with
        vehemence:}
}
{
}

\Verse{66}
{
    \zA{探春早已料定没有自己联的了，便早写出来，因说：“还没 收住呢。”}
}
{
    \eA{No wind doth blow, but yet there is a rush. Pao-ch'in promptly also smiled,
        and strung on: No rain lo falls, but still a patter's heard. Hsiang-yün was
        leaning over,}
}
{
}

\Verse{67}
{
    \zA{李纹听了，接过来，便联了一句道： 欲志今朝乐， 李绮收了一句道： 凭诗祝舜尧。}
}
{
    \eA{indulging in such merriment that she was quite doubled up in two. But
        everybody else had realised that the trio was struggling for mastery, so
        without attempting to versify they kept their gaze fixed on them and gave
        way to laughter.}
}
{
}

\Verse{68}
{
    \zA{李纨道：“够了够了。}
}
{
    \eA{Tai-yü gave her another push to try and induce her to go on.}
}
{
}

\Verse{69}
{
    \zA{虽没作完了韵，腾挪的字，若生扭了，倒不好了。”}
}
{
    \eA{"Do you also sometimes come to your wits' ends; and run to the end of your
        tether?" she went on to say. "I'd like to see what other stuff and nonsense
        you can come out with!"}
}
{
}

\Verse{70}
{
    \zA{说 着大家来细细评论一回，独湘云的多，都笑道：“这都是那块鹿肉的功劳。”}
}
{
    \eA{Hsiang-yün however simply fell forward on Pao-ch'ai's lap and laughed
        incessantly. "If you've got any gumption about you," Pao-ch'ai exclaimed,
        shoving her up,}
}
{
}

\Verse{71}
{
    \zA{李纨 笑道：“逐句评去，却还一气，只是宝玉又落了第了。”}
}
{
    \eA{"take the second rhymes under 'Hsiao' and exhaust them all, and I'll then
        bend the knee to you." "It isn't as if I were writing verses," Hsiang-yün
        laughed rising to her feet;}
}
{
}

\Verse{72}
{
    \zA{宝玉笑道：“我原不会联 句，只好担待我罢。”}
}
{
    \eA{"it's really as if I were fighting for very life." "It's for you to come out
        with something," they all cried with a laugh.}
}
{
}

\Verse{73}
{
    \zA{李纨笑道：“也没有社社担待的：又说‘韵险’了，又整误 了，又‘不会联句’!今日必罚你。}
}
{
    \eA{T'an Ch'un had long ago determined in her mind that there could be no other
        antithetical sentences that she herself could possibly propose, and she
        forthwith set to work to copy out the verses. But as she passed the remark:}
}
{
}

\Verse{74}
{
    \zA{我才看见栊翠庵的红梅有趣，我要折一枝插在 瓶。}
}
{
    \eA{"They haven't as yet been brought to a proper close," Li Wen took up the
        clue, as soon as she caught her words, and added the sentiment:}
}
{
}

\Verse{75}
{
    \zA{可厌妙玉为人，我不理他，如今罚你取一枝来插着玩儿。”}
}
{
    \eA{My wish is to record this morning's fun. Li Ch'i then suggested as a finale
        the line: By these verses, I'd fain sing th' Emperor's praise. "That's
        enough, that will do!"}
}
{
}

\Verse{76}
{
    \zA{众人都道：“这罚 的又雅又有趣。”}
}
{
    \eA{Li Wan cried. "The rhymes haven't, I admit, been exhausted, but any outside
        words you might introduce,}
}
{
}

\Verse{77}
{
    \zA{宝玉也乐为，答应着就要走。}
}
{
    \eA{will, if used in a forced sense, be worth nothing at all." While continuing
        their arguments,}
}
{
}

\Verse{78}
{
    \zA{湘云黛玉一起说着：“外头冷得很， 你且吃杯热酒再去。”}
}
{
    \eA{the various inmates drew near and kept up a searching criticism for a time.
        Hsiang-yün was found to be the one among them, who had devised the largest
        number of lines. "This is mainly due," they unanimously laughed,}
}
{
}

\Verse{79}
{
    \zA{于是湘云早热起壶酒来了，黛玉递了个大杯，满斟了一杯。}
}
{
    \eA{"to the virtue of that piece of venison!" "Let's review them line by line as
        they come," Li Wan smilingly proposed, "but yet as if they formed one
        continuous poem.}
}
{
}

\Verse{80}
{
    \zA{湘云笑道：“你吃了我们这酒，要取不来，加倍罚你！”}
}
{
    \eA{Here's Pao-yü last again!" "I haven't, the fact is, the knack of pairing
        sentences," Pao-yü rejoined with a smile. "You'd better therefore make some
        allowance for me!"}
}
{
}

\Verse{81}
{
    \zA{宝玉忙吃了一杯，冒雪而 去。}
}
{
    \eA{"There's no such thing as making allowances for you in meeting after
        meeting," Li Wan demurred laughing,}
}
{
}

\Verse{82}
{
    \zA{李纨命人好好跟着，黛玉忙拦说：“不必，有了人反不得了。”}
}
{
    \eA{"that you should again after that give out the rhymes in a reckless manner,
        waste your time and not show yourself able to put two lines together. You
        must absolutely bear a penalty today.}
}
{
}

\Verse{83}
{
    \zA{李纨点头道是， 一面命丫鬟将一个美女耸肩瓶拿来，贮了水准备插梅。}
}
{
    \eA{I just caught a glimpse of the red plum in the Lung Ts'ui monastery; and how
        charming it is! I meant to have plucked a twig to put in a vase, but so
        loathsome is the way in which Miao Yü goes on, that I won't have anything to
        do with her!}
}
{
}

\Verse{84}
{
    \zA{因又笑道：“回来该吟红梅 了。”}
}
{
    \eA{But we'll punish him by making him, for the sake of fun, fetch a twig for us
        to put in water."}
}
{
}

\Verse{85}
{
    \zA{湘云忙道：“我先作一首。”}
}
{
    \eA{"This penalty," they shouted with one accord, "is both excellent as well as
        pleasant."}
}
{
}

\Verse{86}
{
    \zA{宝钗笑道：“今日断不容你再作了，你都抢了 去，别人都闲着也没趣。}
}
{
    \eA{Pao-yü himself was no less delighted to carry it into execution, so
        signifying his readiness to comply with their wishes, he felt desirous to be
        off at once. "It's exceedingly cold outside,"}
}
{
}

\Verse{87}
{
    \zA{回来罚宝玉。}
}
{
    \eA{Hsiang-yün and Tai-yü simultaneously remarked,}
}
{
}

\Verse{88}
{
    \zA{他说不会联句，如今就叫他自己做去。”}
}
{
    \eA{"so have a glass of warm wine before you go." Hsiang-yün speedily took up
        the kettle, and Tai-yü handed him a large cup,}
}
{
}

\Verse{89}
{
    \zA{黛 玉笑道：“这话很是。}
}
{
    \eA{filled to the very brim. "Now swallow the wine we give you,"}
}
{
}

\Verse{90}
{
    \zA{我还有主意：方才联句不够，莫若拣那联得少的人作红梅诗。”}
}
{
    \eA{Hsiang-yün smiled. "And if you don't bring any plum blossom, we'll inflict a
        double penalty." Pao-yü gulped down hurry-scurry the whole contents of the
        cup and started on his errand in the face of the snow.}
}
{
}

\Verse{91}
{
    \zA{宝钗笑道：“这话是极。}
}
{
    \eA{"Follow him carefully." Li Wan enjoined the servants.}
}
{
}

\Verse{92}
{
    \zA{方才邢李三位屈才，且又是客，琴儿和颦儿云儿抢了他们 许多。}
}
{
    \eA{Tai-yü, however, hastened to interfere and make her desist. "There's no such
        need," she cried. "Were any one to go with him, he'll contrariwise not get
        the flowers." Li Wan nodded her head. "Yes!"}
}
{
}

\Verse{93}
{
    \zA{我们一概都别作，只他们三人做才是。”}
}
{
    \eA{she assented, and then went on to direct a waiting-maid to bring a vase, in
        the shape of a beautiful girl with high shoulders,}
}
{
}

\Verse{94}
{
    \zA{李纨因说：“绮儿也不大会做，还 是让琴妹妹罢。”}
}
{
    \eA{to fill it with water, and get it ready to put the plum blossom in. "And
        when he comes back," she felt induced to add, "we must recite verses on the
        red plum."}
}
{
}

\Verse{95}
{
    \zA{宝钗只得依允。}
}
{
    \eA{"I'll indite a stanza in advance," eagerly exclaimed Hsiang-yün.}
}
{
}

\Verse{96}
{
    \zA{又道：“就用‘红梅花’三个字做韵，每人一首 七言律：邢大妹妹做‘红’字，你们李大妹妹做‘梅’字，琴儿做‘花’字。”}
}
{
    \eA{"We'll on no account let you indite any more to-day," Pao-ch'ai laughed.
        "You beat every one of us hollow; so if we sit with idle hands, there won't
        be any fun. But by and bye we'll fine Pao-yü; and, as he says that he can't
        pair antithetical lines,}
}
{
}

\Verse{97}
{
    \zA{李 纨道：“饶过宝玉去，我不服。”}
}
{
    \eA{we'll now make him compose a stanza himself." "This is a capital idea!"
        Tai-yü smiled. "But I've got another proposal.}
}
{
}

\Verse{98}
{
    \zA{湘云忙道：“有个好题目命他做。”}
}
{
    \eA{As the lines just paired are not sufficient, won't it be well to pick out
        those who've put together the fewest distiches,}
}
{
}

\Verse{99}
{
    \zA{众人问：“何 题？”}
}
{
    \eA{and make them versify on the red plum blossom?"}
}
{
}

\Verse{100}
{
    \zA{湘云道；“命他就做‘访妙玉乞红梅’，岂不有趣？”}
}
{
    \eA{"An excellent proposal!" Pao-ch'ai ventured laughing. "The three girls Hsing
        Chou-yen, Li Wen and Li Ch'i, failed just now to do justice to their
        talents;}
}
{
}

\Verse{101}
{
    \zA{众人听了，都说：“有 趣！”}
}
{
    \eA{besides they are visitors; and as Ch'in Erh, P'in Erh and Yün Erh got the
        best of us by a good deal,}
}
{
}

\Verse{102}
{
    \zA{一语未了，只见宝玉笑欣欣擎了一枝红梅进来。}
}
{
    \eA{it's only right that none of us should compose any more, and that that trio
        should only do so." "Ch'i Erh," Li Wan thereupon retorted,}
}
{
}

\Verse{103}
{
    \zA{众丫鬟忙已接过，插入瓶内。}
}
{
    \eA{"is also not a very good hand at verses, let therefore cousin Ch'in have a
        try!"}
}
{
}
\ChapterImage{img/chapters/109第50回 琉璃世界白雪紅梅 脂粉香娃割腥啖膻.jpg}


\Verse{104}
{
    \zA{众人都道：“来赏玩！”}
}
{
    \eA{Pao-ch'ai had no alternative but to express her acquiescence.}
}
{
}

\Verse{105}
{
    \zA{宝玉笑道：“你们如今赏罢，也不知费了我多少精神呢。”}
}
{
    \eA{"Let the three words 'red plum blossom,'" she then suggested, "be used for
        rhymes; and let each person compose an heptameter stanza. Cousin Hsing to
        indite on the word 'red;'}
}
{
}

\Verse{106}
{
    \zA{说着，探春早又递了一钟暖酒来，众丫鬟上来接了蓑笠掸雪。}
}
{
    \eA{your elder cousin Li on 'plum;' and Ch'in Erh on 'blossom.'" "If you let
        Pao-yü off," Li Wan interposed, "I won't have it!" "I've got a capital
        theme," Hsiung-yün eagerly remarked,}
}
{
}

\Verse{107}
{
    \zA{各人屋里丫鬟都添送 衣裳来，袭人也遣人送了半旧的狐腋褂来。}
}
{
    \eA{"so let's make him write some!" "What theme is it?" one and all inquired.
        "If we made him," Hsiang-yün resumed, "versify on: 'In search of Miao Yü to
        beg for red plum blossom,'}
}
{
}

\Verse{108}
{
    \zA{李纨命人将那蒸的大芋头盛了一盘，又 将朱桔、黄橙、橄榄等物盛了两盘，命人带给袭人去。}
}
{
    \eA{won't it be full of fun?" "That will be full of zest," the party exclaimed,
        upon hearing the theme propounded by her. But hardly had they given
        expression to their approval than they perceived Pao-yü come in, beaming
        with smiles and glee,}
}
{
}

\Verse{109}
{
    \zA{湘云且告诉宝玉方才的诗题， 又催宝玉快做。}
}
{
    \eA{and holding with both hands a branch of red plum blossom. The maids
        hurriedly relieved him of his burden and put the branch in the vase,}
}
{
}

\Verse{110}
{
    \zA{宝玉道：“好姐姐好妹妹们，让我自己用韵罢，别限韵了。”}
}
{
    \eA{and the inmates present came over in a body to feast their eyes on it.
        "Well, may you look at it now," Pao-yü smiled. "You've no idea what an
        amount of trouble it has cost me!"}
}
{
}

\Verse{111}
{
    \zA{众人 都说：“随你做去罢。”}
}
{
    \eA{As he uttered these words, T'an Ch'un handed him at once another cup of warm
        wine;}
}
{
}

\Verse{112}
{
    \zA{一面说，一面大家看梅花。}
}
{
    \eA{and the maids approached, and took his wrapper and hat,}
}
{
}

\Verse{113}
{
    \zA{原来这一枝梅花只有二尺来高， 旁有一枝纵横而出，约有二三尺长，其间小枝分歧，或如蟠螭，或如僵蚓，或孤削 如笔，或密聚如林，真乃花吐胭脂，香欺兰蕙。}
}
{
    \eA{and shook off the snow. But the servant-girls attached to their respective
        quarters then brought them over extra articles of clothing. Hsi Jen, in like
        manner, despatched a domestic with a pelisse, the worse for wear, lined with
        fur from foxes' ribs, so Li Wan, having directed a servant to fill a plate
        with steamed large taros, and to make up two dishes with red-skinned
        oranges, yellow coolie oranges, olives and other like things,}
}
{
}

\Verse{114}
{
    \zA{各各称赏。}
}
{
    \eA{bade some one take them over to Hsi Jen.}
}
{
}

\Verse{115}
{
    \zA{谁知岫烟、李纹、宝琴三人都已吟成，各自写了出来。}
}
{
    \eA{Hsiang-yün also communicated to Pao-yü the subject for verses they had
        decided upon a short while back. But she likewise urged Pao-yü to be quick
        and accomplish his task.}
}
{
}

\Verse{116}
{
    \zA{众人便依“红”“梅” “花”三字之序看去，写道： 赋得红梅花　　邢岫烟 桃未芳菲杏未红，冲寒先喜笑东风。}
}
{
    \eA{"Dear senior cousin, dear junior cousin," pleaded Pao-yü, "let me use my own
        rhymes. Don't bind me down to any." "Go on as you like," they replied with
        one consent. But conversing the while, they passed the plum blossom under
        inspection.}
}
{
}

\Verse{117}
{
    \zA{魂飞庾岭春难辨，霞隔罗浮梦未通。}
}
{
    \eA{This bough of plum blossom was, in fact, only two feet in height; but from
        the side projected a branch,}
}
{
}

\Verse{118}
{
    \zA{绿萼添妆融宝炬，缟仙扶醉跨残虹。}
}
{
    \eA{crosswise, about two or three feet in length the small twigs and stalks on
        which resembled coiled dragons,}
}
{
}

\Verse{119}
{
    \zA{看来岂是寻常色，浓淡由他冰雪中。}
}
{
    \eA{or crouching earthworms; and were either single and trimmed pencil-like, or
        thick and bushy grove-like.}
}
{
}

\Verse{120}
{
    \zA{又　　李　纹 白梅懒赋赋红梅，逞艳先迎醉眼开。}
}
{
    \eA{Indeed, their appearance was as if the blossom spurted cosmetic. This
        fragrance put orchids to the blush.}
}
{
}

\Verse{121}
{
    \zA{冻脸有痕皆是血，酸心无恨亦成灰。}
}
{
    \eA{So every one present contributed her quota of praise. Chou-yen, Li Wen and
        Pao-ch'in had,}
}
{
}

\Verse{122}
{
    \zA{误吞丹药移真骨，偷下瑶池脱旧胎。}
}
{
    \eA{little though it was expected, all three already finished their lines and
        each copied them out for herself,}
}
{
}

\Verse{123}
{
    \zA{江北江南春灿烂，寄言蜂蝶漫疑猜。}
}
{
    \eA{so the company began to peruse their compositions, subjoined below, in the
        order of the three words:}
}
{
}

\Verse{124}
{
    \zA{又　　宝　琴 疏是枝条艳是花，春妆儿女竞奢华。}
}
{
    \eA{'red plum blossom.' Verses to the red plum blossom by Hsing Chou-yen. The
        peach tree has not donned its fragrance yet, the almond is not red. What
        time it strikes the cold,}
}
{
}

\Verse{125}
{
    \zA{闲庭曲槛无馀雪，流水空山有落霞。}
}
{
    \eA{it's first joyful to smile at the east wind. When its spirit to the Yü Ling
        hath flown, 'tis hard to say 'tis spring.}
}
{
}

\Verse{126}
{
    \zA{幽梦冷随红袖笛，游仙香泛绛河槎。}
}
{
    \eA{The russet clouds across the 'Lo Fu' lie, so e'en to dreams it's closed. The
        green petals add grace to a coiffure,}
}
{
}

\Verse{127}
{
    \zA{前身定是瑶台种，无复相疑色相差。}
}
{
    \eA{when painted candles burn. The simple elf when primed with wine doth the
        waning rainbow bestride.}
}
{
}

\Verse{128}
{
    \zA{众人看了，都笑着称赞了一回，又指末一首更好。}
}
{
    \eA{Does its appearance speak of a colour of ordinary run? Both dark and light
        fall of their own free will into the ice and snow. The next was the
        production of Li Wen,}
}
{
}

\Verse{129}
{
    \zA{宝玉见宝琴年纪最小，才又 敏捷；黛玉湘云二个斟了一小杯酒，都贺宝琴。}
}
{
    \eA{and its burden was: To write on the white plum I'm not disposed, but I'll
        write on the red. Proud of its beauteous charms, 'tis first to meet the
        opening drunken eye. On its frost-nipped face are marks;}
}
{
}

\Verse{130}
{
    \zA{宝钗笑道：“三首各有好处。}
}
{
    \eA{and these consist wholly of blood. Its heart is sore, but no anger it knows;}
}
{
}

\Verse{131}
{
    \zA{你们 两个天天捉弄厌了我，如今又捉弄他来了。”}
}
{
    \eA{to ashes too it turns. By some mistake a pill (a fairy) takes and quits her
        real frame. From the fairyland pool she secret drops, and casts off her old
        form.}
}
{
}

\Verse{132}
{
    \zA{李纨又问宝玉：“你可有了？”}
}
{
    \eA{In spring, both north and south of the river, with splendour it doth bloom.}
}
{
}

\Verse{133}
{
    \zA{宝玉忙道；“我倒有了，才一看见这三首，又 唬忘了。}
}
{
    \eA{Send word to bees and butterflies that they need not give way to fears! This
        stanza came next from the pen of Hsüeh Pao-ch'in, Far distant do the
        branches grow;}
}
{
}

\Verse{134}
{
    \zA{等我再想。”}
}
{
    \eA{but how beauteous the blossom blooms!}
}
{
}

\Verse{135}
{
    \zA{湘云听了，便拿了一支铜火箸击着手炉，笑道：“我击了， 若鼓绝不成，又要罚的。”}
}
{
    \eA{The maidens try with profuse show to compete in their spring head-dress. No
        snow remains on the vacant pavilion and the tortuous rails. Upon the running
        stream and desolate hills descend the russet clouds. When cold prevails one
        can in a still dream follow the lass-blown fife.}
}
{
}

\Verse{136}
{
    \zA{宝玉笑道：“我已有了。}
}
{
    \eA{The wandering elf roweth in fragrant spring, the boat in the red stream.}
}
{
}

\Verse{137}
{
    \zA{“黛玉提起笔来，笑道：“你 念我写。”}
}
{
    \eA{In a previous existence, it must sure have been of fairy form. No doubt need
        'gain arise as to its beauty differing from then.}
}
{
}

\Verse{138}
{
    \zA{湘云便击了一下，笑道：“一鼓绝。”}
}
{
    \eA{The perusal over, they spent some time in heaping, smiling the while,
        eulogiums upon the compositions.}
}
{
}

\Verse{139}
{
    \zA{宝玉笑道：“有了，你写罢。”}
}
{
    \eA{And they pointed at the last stanza as the best of the lot; which made it
        evident to Pao-yü that Pao-ch'in,}
}
{
}

\Verse{140}
{
    \zA{众人听他念道： 酒未开樽句未裁， 黛玉写了，摇头笑道：“起的平平。”}
}
{
    \eA{albeit the youngest in years, was, on the other hand, the quickest in wits.
        Tai-yü and Hsiang-yün then filled up a small cup with wine and
        simultaneously offered their congratulations to Pao-ch'in.}
}
{
}

\Verse{141}
{
    \zA{湘云又道：“快着。”}
}
{
    \eA{"Each of the three stanzas has its beauty,"}
}
{
}

\Verse{142}
{
    \zA{宝玉笑道： 寻春问腊到蓬莱。}
}
{
    \eA{Pao-ch'ai remarked, a smile playing round her lips. "You two have daily made
        a fool of me,}
}
{
}

\Verse{143}
{
    \zA{黛玉湘云都点头笑道：“有些意思了。”}
}
{
    \eA{and are you now going to fool her also?" "Have you got yours ready?" Li Wan
        went on to inquire of Pao-yü.}
}
{
}

\Verse{144}
{
    \zA{宝玉又道： 不求大士瓶中露，为乞孀娥槛外梅。}
}
{
    \eA{"I'd got them," Pao-yü promptly answered, "but the moment I read their three
        stanzas, I once more became so nervous that they quite slipped from my mind.}
}
{
}

\Verse{145}
{
    \zA{黛玉写了，摇头说：“小巧而已。”}
}
{
    \eA{But let me think again." Hsiang-yün, at this reply, fetched a copper poker,
        and,}
}
{
}

\Verse{146}
{
    \zA{湘云将手又敲了一下。}
}
{
    \eA{while beating on the hand-stove, she laughingly said: "I shall go on
        tattooing.}
}
{
}

\Verse{147}
{
    \zA{宝玉笑道： 入世冷挑红雪去，离尘香割紫云来。}
}
{
    \eA{Now mind if when the drumming ceases, you haven't accomplished your task,
        you'll have to bear another fine."}
}
{
}

\Verse{148}
{
    \zA{槎？}
}
{
    \eA{"I've already got them!"}
}
{
}

\Verse{149}
{
    \zA{谁惜诗肩瘦，衣上犹沾佛院苔。}
}
{
    \eA{Pao-yü rejoined, smilingly. Tai-yü then picked up a pencil. "Recite them,"
        she smiled,}
}
{
}

\Verse{150}
{
    \zA{黛玉写毕，湘云大家才评论时，只见几个丫鬟跑进来道：“老太太来了。”}
}
{
    \eA{"and I'll write them down." Hsiang-yün beat one stroke (on the stove). "The
        first tattoo is over," she laughed. "I'm ready," Pao-yü smiled. "Go on
        writing."}
}
{
}

\Verse{151}
{
    \zA{众 人忙迎出来，大家又笑道：“怎么这等高兴！”}
}
{
    \eA{At this, they heard him recite: The wine bottle is not opened, the line is
        not put into shape. Tai-yü noted it down, and shaking her head,}
}
{
}

\Verse{152}
{
    \zA{说着，远远见贾母围了大斗篷，带 着灰鼠暖兜，坐着小竹轿，打着青绸油伞，鸳鸯琥珀等五六个丫鬟，每人都是打着 伞，拥轿而来。}
}
{
    \eA{"They begin very smoothly," she said, as she smiled. "Be quick!" Hsiang-yün
        again urged. Pao-yü laughingly continued: To fairyland I speed to seek for
        spring,}
}
{
}

\Verse{153}
{
    \zA{李纨等忙往上迎。}
}
{
    \eA{and the twelfth moon to find. Tai-yü and Hsiang-yün both nodded.}
}
{
}

\Verse{154}
{
    \zA{贾母命人止住，说：“只站在那里就是了。”}
}
{
    \eA{"It's rather good," they smiled. Pao-yü resumed, saying: I will not beg the
        high god for a bottle of the (healing) dew,}
}
{
}

\Verse{155}
{
    \zA{来 至跟前，贾母笑道：“我瞒着你太太和凤丫头来了。}
}
{
    \eA{But pray Shuang O to give me some plum bloom beyond the rails. Tai-yü jotted
        the lines down and wagged her head to and fro. "They're ingenious, that's
        all,"}
}
{
}

\Verse{156}
{
    \zA{大雪地下，我坐着这个无妨， 没的叫他娘儿们踩雪吗。”}
}
{
    \eA{she observed. Hsiang-yün gave another rap with her hand. Pao-yü thereupon
        smilingly added: I come into the world and, in the cold, I pick out some red
        snow.}
}
{
}

\Verse{157}
{
    \zA{众人忙上前来接斗篷，搀扶着，一面答应着。}
}
{
    \eA{I leave the dusty sphere and speed to pluck the fragrant purple clouds. I
        bring a jagged branch, but who in pity sings my shoulders thin?}
}
{
}

\Verse{158}
{
    \zA{贾母来至室中，先笑道：“好俊梅花!你们也会乐，我也不饶你们！”}
}
{
    \eA{On my clothes still sticketh the moss from yon Buddhistic court. As soon as
        Tai-yü had done writing, Hsiang-yün and the rest of the company began to
        discuss the merits of the verses;}
}
{
}

\Verse{159}
{
    \zA{说着， 李纨早命人拿了一个大狼皮褥子来，铺在当中。}
}
{
    \eA{but they then saw several servant-maids rush in, shouting: "Our venerable
        mistress has come." One and all hurried out with all despatch to meet her.}
}
{
}

\Verse{160}
{
    \zA{贾母坐了，因笑道：“你们只管照 旧玩笑吃喝。}
}
{
    \eA{"How comes it that she is in such good cheer?" every one also laughed.
        Speaking the while, they discerned, at a great distance,}
}
{
}

\Verse{161}
{
    \zA{我因为天短了，不敢睡中觉，抹了一会牌，想起你们来了，我也来凑 个趣儿。”}
}
{
    \eA{their grandmother Chia seated, enveloped in a capacious wrapper, and rolled
        up in a warm hood lined with squirrel fur, in a small bamboo sedan-chair
        with an open green silk glazed umbrella in her hand.}
}
{
}

\Verse{162}
{
    \zA{李纨早又捧过手炉来。}
}
{
    \eA{Yüan Yang, Hu Po and some other girls, mustering in all five or six,}
}
{
}

\Verse{163}
{
    \zA{探春另拿了一副杯箸来，亲自斟了暖酒奉给贾母。}
}
{
    \eA{held each an umbrella and pressed round the chair, as they advanced. Li Wan
        and her companions went up to them with hasty step; but dowager lady Chia
        directed the servants to make them stop;}
}
{
}

\Verse{164}
{
    \zA{贾母便饮了一口，问：“那个盘子是什么东西？”}
}
{
    \eA{explaining that it would be quite enough if they stood where they were. On
        her approach, old lady Chia smiled. "I've given," she observed,}
}
{
}

\Verse{165}
{
    \zA{众人忙捧了过来回说：“是糟鹌 鹑。”}
}
{
    \eA{"your Madame Wang and that girl Feng the slip and come. What deep snow
        covers the ground! For me, I'm seated in this,}
}
{
}

\Verse{166}
{
    \zA{贾母道：“这倒罢了，撕一点子腿儿来。”}
}
{
    \eA{so it doesn't matter; but you mustn't let those ladies trudge in the snow."
        The various followers rushed forward to take her wrapper and to support her,}
}
{
}

\Verse{167}
{
    \zA{李纨忙答应了，要水洗手，亲自 来撕。}
}
{
    \eA{and as they did so, they expressed their acquiescence. As soon as she got
        indoors old lady Chia was the first to exclaim with a beaming face:}
}
{
}

\Verse{168}
{
    \zA{贾母道：“你们仍旧坐下说笑，我听着才喜欢。”}
}
{
    \eA{"What beautiful plum blossom! You well know how to make merry; but I too
        won't let you off!" But in the course of her remarks,}
}
{
}

\Verse{169}
{
    \zA{又命李纨：“你也只管坐 下，就如同我没来的一样才好，不然我就走了。”}
}
{
    \eA{Li Wan quickly gave orders to a domestic to fetch a large wolf skin rug, and
        to spread it in the centre, so dowager lady Chia made herself comfortable on
        it. "Just go on as before with your romping and joking,}
}
{
}

\Verse{170}
{
    \zA{众人听了，方才依次坐下，只李 纨挪到尽下边。}
}
{
    \eA{drinking and eating," she then laughed. "As the days are so short, I did not
        venture to have a midday siesta. After therefore playing at dominoes for a
        time,}
}
{
}

\Verse{171}
{
    \zA{贾母因问：“你们作什么玩呢？”}
}
{
    \eA{I bethought myself of you people, and likewise came to join the fun." Li Wan
        soon also presented her a hand-stove,}
}
{
}

\Verse{172}
{
    \zA{众人便说：“做诗呢。”}
}
{
    \eA{while T'an Ch'un brought an extra set of cups and chopsticks,}
}
{
}

\Verse{173}
{
    \zA{贾母道： “有做诗的，不如做些灯谜儿，大家正月里好玩。”}
}
{
    \eA{and filling with her own hands, a cup with warm wine, she handed it to her
        grandmother Chia. Old lady Chia swallowed a sip. "What's there in that
        dish?"}
}
{
}

\Verse{174}
{
    \zA{众人答应。}
}
{
    \eA{she afterwards inquired.}
}
{
}

\Verse{175}
{
    \zA{说笑了一会，贾母 便说：“这里潮湿，你们别久坐，仔细着了凉。}
}
{
    \eA{The various inmates hurriedly carried it over to her, and explained that
        'they were pickled quails.' "These won't hurt me," dowager lady Chia said,
        "so cut off a piece of the leg and give it to me."}
}
{
}

\Verse{176}
{
    \zA{倒是你四妹妹那里暖和，我们到那 里瞧瞧他的画儿，赶年可能有了不能。”}
}
{
    \eA{"Yes!" promptly acquiesced Li Wan, and asking for water, she washed her
        hands, and then came in person to carve the quail. "Sit down again," dowager
        lady Chia said, pressing them, "and go on with your chatting and laughing.}
}
{
}

\Verse{177}
{
    \zA{众人笑道：“那里能年下就有了？}
}
{
    \eA{Let me hear you, and feel happy. Just you also seat yourself," continuing,
        she remarked to Li Wan,}
}
{
}

\Verse{178}
{
    \zA{只怕明 年端阳才有呢。”}
}
{
    \eA{"and behave as if I were not here. If you do so, well and good. Otherwise, I
        shall take myself off at once."}
}
{
}

\Verse{179}
{
    \zA{贾母道：“这还了得，他竟比盖这园子还费工夫了。”}
}
{
    \eA{But it was only when they heard how persistent she was in her solicitations
        that they all resumed the seats, which accorded with their age, with the
        exception of Li Wan, who moved to the furthest side.}
}
{
}

\Verse{180}
{
    \zA{说着，仍坐了竹椅轿，大家围随，过了藕香榭，穿入一条夹道，东西两边皆是 过街门，门楼上里外都嵌着石头匾。}
}
{
    \eA{"What were you playing at?" old lady Chia thereupon asked. "We were writing
        verses," answered the whole party. "Wouldn't it be well for those who are up
        to poetry," dowager lady Chia suggested; "to devise a few puns for lanterns
        so that the whole lot of us should be able to have some fun in the first
        moon?"}
}
{
}

\Verse{181}
{
    \zA{如今进的是西门，向外的匾上凿着“穿云”二 字，向里的凿着“度月”两字。}
}
{
    \eA{With one voice, they expressed their approval. But after they had jested for
        a little time; "It's damp in here;" old lady Chia said, "so don't you sit
        long, for mind you might be catching cold.}
}
{
}

\Verse{182}
{
    \zA{来至堂中，进了向南的正门，贾母下了轿，惜春已 接出来了。}
}
{
    \eA{Where it's nice and warm is in your cousin Quarta's over there, so let's all
        go and see how she is getting on with her painting, and whether it will be
        ready or not by the end of the year."}
}
{
}

\Verse{183}
{
    \zA{从里面游廊过去，便是惜春卧房，厦檐下挂着“暖香坞”的匾，早有几 个人打起猩红毡帘，已觉暖气拂脸。}
}
{
    \eA{"How could it be completed by the close of the year?" they smiled. "She
        could only, we fancy, get it ready by the dragon boat festival next year."
        "This is dreadful!" old lady Chia exclaimed. "Why, she has really wasted
        more labour on it than would have been actually required to lay out this
        garden!"}
}
{
}

\Verse{184}
{
    \zA{大家进入屋里，贾母并不归坐，只问惜春：“画 到那里了？”}
}
{
    \eA{With these words still on her lips, she ensconced herself again in the
        bamboo sedan, and closed in or followed by the whole company, she repaired
        to the Lotus Fragrance Arbour,}
}
{
}

\Verse{185}
{
    \zA{惜春因笑回：“天气寒冷了，胶性都凝涩不润，画了恐不好看，故此 收起来了。”}
}
{
    \eA{where they got into a narrow passage, flanked on the east as well as the
        west, with doors from which they could cross the street. Over these doorways
        on the inside as well as outside were inserted alike tablets made of stone.}
}
{
}

\Verse{186}
{
    \zA{贾母笑道：“我年下就要的，你别脱懒儿，快拿出来给我快画。”}
}
{
    \eA{The door they went in by, on this occasion, lay on the west. On the tablet
        facing outwards, were cut out the two words representing: 'Penetrating into
        the clouds.'}
}
{
}

\Verse{187}
{
    \zA{一 语未了，忽见凤姐披着紫羯绒褂笑嘻嘻的来了，口内说道：“老祖宗今儿也不告诉 人，私自就来了，叫我好找！”}
}
{
    \eA{On that inside, were engraved the two characters meaning: 'crossing to the
        moon.' On their arrival at the hall, they walked in by the main entrance,
        which looked towards the south. Dowager lady Chia then alighted from her
        chair. Hsi Ch'un had already made her appearance out of doors to welcome
        her,}
}
{
}

\Verse{188}
{
    \zA{贾母见他来了，心中喜欢，道：“我怕你冻着，所 以不许人告诉你去。}
}
{
    \eA{so taking the inner covered passage, they passed over to the other side and
        reached Hsi Ch'un's bedroom; on the door posts of which figured the three
        words: 'Warm fragrance isle.' Several servants were at once at hand;}
}
{
}

\Verse{189}
{
    \zA{你真是个小鬼灵精儿，到底找了我来。}
}
{
    \eA{and no sooner had they raised the red woollen portière, than a soft
        fragrance wafted itself into their faces.}
}
{
}

\Verse{190}
{
    \zA{论礼，孝敬也不在这上 头。”}
}
{
    \eA{The various inmates stepped into the room. Old lady Chia, however, did not
        take a seat,}
}
{
}

\Verse{191}
{
    \zA{凤姐儿笑道：“我那里是孝敬的心找了来呢？}
}
{
    \eA{but simply inquired where the painting was. "The weather is so bitterly
        cold," Hsi Ch'un consequently explained smiling,}
}
{
}

\Verse{192}
{
    \zA{我因为到了老祖宗那里，鸦没 雀静的，问小丫头子们，他又不肯叫我找到园里来。}
}
{
    \eA{"that the glue, whose property is mainly to coagulate, cannot be moistened,
        so I feared that, were I to have gone on with the painting, it wouldn't be
        worth looking at; and I therefore put it away." "I must have it by the close
        of the year,"}
}
{
}

\Verse{193}
{
    \zA{我正疑惑，忽然又来了两个姑 子。}
}
{
    \eA{dowager lady Chia laughed, "so don't idle your time away. Produce it at once
        and go on painting for me,}
}
{
}

\Verse{194}
{
    \zA{我心里才明白了，那姑子必是来送年疏或要年例香例银子，老祖宗年下的事也 多，一定是躲债来了。}
}
{
    \eA{as quick as you can." But scarcely had she concluded her remark, than she
        unexpectedly perceived lady Feng arrive, smirking and laughing, with a
        purple pelisse, lined with deer fur, thrown over her shoulders. "Venerable
        senior!" she shouted, "You don't even so much as let any one know to-day,}
}
{
}

\Verse{195}
{
    \zA{我赶忙问了那姑子，果然不错。}
}
{
    \eA{but sneak over stealthily. I've had a good hunt for you!" When old lady Chia
        saw her join them,}
}
{
}

\Verse{196}
{
    \zA{我才就把年例给了他们去了。}
}
{
    \eA{she felt filled with delight. "I was afraid," she rejoined, "that you'd be
        feeling cold.}
}
{
}

\Verse{197}
{
    \zA{这会子老祖宗的债主儿已去了，不用躲着了。}
}
{
    \eA{That's why, I didn't allow any one to tell you. You're really as sharp as a
        spirit to have, at last, been able to trace my whereabouts!}
}
{
}

\Verse{198}
{
    \zA{已预备下稀嫩的野鸡，请用晚饭去罢， 再迟一回就老了。”}
}
{
    \eA{But according to strict etiquette, you shouldn't show filial piety to such a
        degree!" "Is it out of any idea of filial piety that I came after you?}
}
{
}

\Verse{199}
{
    \zA{他一行说，众人一行笑。}
}
{
    \eA{Not at all!" lady Feng added with a laugh. "But when I got to your place,}
}
{
}

\Verse{200}
{
    \zA{凤姐儿也不等贾母说话，便命人抬过轿来。}
}
{
    \eA{worthy senior, I found everything so quiet that not even the caw of a crow
        could be heard, and when I asked the young maids where you'd gone,}
}
{
}

\Verse{201}
{
    \zA{贾母笑着 挽了凤姐儿的手，仍上了轿，带着众人，说笑出了夹道东门。}
}
{
    \eA{they wouldn't let me come and search in the garden. So I began to give way
        to surmises. Suddenly also arrived two or three nuns; and then, at length, I
        jumped at the conclusion that these women must have come to bring their
        yearly prayers,}
}
{
}

\Verse{202}
{
    \zA{一看四面，粉妆银砌， 忽见宝琴披着凫靥裘，站在山坡背后遥等，身后一个丫鬟，抱着一瓶红梅。}
}
{
    \eA{or to ask for their annual or incense allowance, and that, with the amount
        of things you also, venerable ancestor, have to do for the end of the year,
        you had for certain got out of the way of your debts. Speedily therefore I
        inquired of the nuns what it was that brought them there,}
}
{
}

\Verse{203}
{
    \zA{众人都 笑道：“怪道少了两个，他却在这里等着，――也弄梅花去了！”}
}
{
    \eA{and, for a fact, there was no mistake in my surmises. So promptly issuing
        the annual allowances to them, I now come to report to you, worthy senior,
        that your creditors have gone,}
}
{
}

\Verse{204}
{
    \zA{贾母喜的忙笑道； “你们瞧，这雪坡儿上，配上他这个人物儿，又是这件衣裳，后头又是这梅花，像 个什么？”}
}
{
    \eA{and that there's no need for you to skulk away. But I've had some tender
        pheasant prepared; so please come, and have your evening meal; for if you
        delay any longer, it will get quite stale."}
}
{
}

\Verse{205}
{
    \zA{众人都笑道：“就像老太太屋里挂的仇十洲画的《艳雪图》。”}
}
{
    \eA{As she spoke, everybody burst out laughing. But lady Feng did not allow any
        time to dowager lady Chia to pass any observations, but forthwith directed
        the servants to bring the chair over.}
}
{
}

\Verse{206}
{
    \zA{贾母摇 头笑道：“那画的那里有这件衣裳？}
}
{
    \eA{Old lady Chia then smilingly laid hold of lady Feng's hand and got again
        into her chair; but she took along with her the whole company of relatives
        for a chat and a laugh.}
}
{
}

\Verse{207}
{
    \zA{人也不能这样好。”}
}
{
    \eA{Upon issuing out of the gate on the east side of the narrow passage,}
}
{
}
\ChapterImage{img/chapters/110第50回 觀景遠望如豔雪圖.jpg}


\Verse{208}
{
    \zA{一语未了，只见宝琴身后 又转出一个穿大红猩猩毡的人来。}
}
{
    \eA{the four quarters presented to their gaze the appearance of being adorned
        with powder, and inlaid with silver. Unawares, they caught sight of
        Pao-ch'in,}
}
{
}

\Verse{209}
{
    \zA{贾母道：“那又是那个女孩儿？”}
}
{
    \eA{in a duck down cloak, waiting at a distance at the back of the hill slope;
        while behind her stood a maid,}
}
{
}

\Verse{210}
{
    \zA{众人笑道：“我 们都在这里，那是宝玉。”}
}
{
    \eA{holding a vase full of red plum blossoms. "Strange enough," they all
        exclaimed laughingly, "two of us were missing!}
}
{
}

\Verse{211}
{
    \zA{贾母笑道：“我的眼越发花了。”}
}
{
    \eA{But she's waiting over there. She's also been after some plum-blossom."
        "Just look," dowager lady Chia eagerly cried out joyfully,}
}
{
}

\Verse{212}
{
    \zA{说话之间，来至跟前， 可不是宝玉和宝琴两个？}
}
{
    \eA{"that human creature has been put there to match with the snow-covered hill!
        But with that costume, and the plum-blossom at the back of her, to what does
        she bear a resemblance?"}
}
{
}

\Verse{213}
{
    \zA{宝玉笑向宝钗黛玉等道：“我才又到了栊翠庵，妙玉竟每 人送你们一枝梅花，我已经打发人送去了。”}
}
{
    \eA{"She resembles," one and all smiled, "Chou Shih-ch'ou's beautiful snow
        picture, suspended in your apartments, venerable ancestor." "Is there in
        that picture any such costume?" Old lady Chia demurred, nodding her head and
        smiling. "What's more the persons represented in it could never be so
        pretty!"}
}
{
}

\Verse{214}
{
    \zA{众人都笑说：“多谢你费心。”}
}
{
    \eA{Hardly had this remark dropped from her mouth, than she discerned some one
        else, clad in a deep red woollen cloak,}
}
{
}

\Verse{215}
{
    \zA{说话之间，已出了园门，来至贾母房中。}
}
{
    \eA{appear to view at the back of Pao-ch'in. "What other girl is that?" dowager
        lady Chia asked. "We girls are all here."}
}
{
}

\Verse{216}
{
    \zA{吃毕饭大家又说笑了一回，忽见薛姨 妈也来了，说：“好大雪，一日也没过来望候老太太。}
}
{
    \eA{they laughingly answered. "That's Pao-yü." "My eyes," old lady Chia smiled,
        "are getting dimmer and dimmer!" So saying, they drew near, and of course,
        they turned out to be Pao-yü and Pao-ch'in. "I've just been again to the
        Lung Ts'ui monastery,"}
}
{
}

\Verse{217}
{
    \zA{今日老太太倒不高兴？}
}
{
    \eA{Pao-yü smiled to Pao-ch'ai, Tai-yü and his other cousins,}
}
{
}

\Verse{218}
{
    \zA{正该 赏雪才是。”}
}
{
    \eA{"and Miao Yü gave me for each of you a twig of plum blossom.}
}
{
}

\Verse{219}
{
    \zA{贾母笑道：“何曾不高兴了!我找了他们姐妹去玩了一会子。”}
}
{
    \eA{I've already sent a servant to take them over." "Many thanks for the trouble
        you've been put to," they, with one voice, replied. But speaking the while,
        they sallied out of the garden gate,}
}
{
}

\Verse{220}
{
    \zA{薛姨 妈笑道：“昨儿晚上我原想着今日要和我们姨太太借一天园子，摆两桌粗酒，请老 太太赏雪的；又见老太太安歇的早，我听见宝儿说：‘老太太心里不大爽。’}
}
{
    \eA{and repaired to their grandmother Chia's suite of apartments. Their meal
        over, they joined in a further chat and laugh, when unexpectedly they saw
        Mrs. Hsüeh also arrive. "With all this snow," she observed, "I haven't been
        over the whole day to see how you, venerable senior, were getting on.}
}
{
}

\Verse{221}
{
    \zA{因此 如今也不敢惊动。}
}
{
    \eA{Your ladyship couldn't have been in a good sort of mood to-day,}
}
{
}

\Verse{222}
{
    \zA{早知如此，我竟该请了才是呢。”}
}
{
    \eA{for you should have gone and seen the snow." "How not in a good mood?" old
        lady Chia exclaimed.}
}
{
}

\Verse{223}
{
    \zA{贾母笑道：“这才是十月，是 头场雪，往后下雪的日子多着呢，再破费姨太太不迟。”}
}
{
    \eA{"I went and looked up these young ladies and had a romp with them for a
        time." "Last night," Mrs. Hsüeh smiled, "I was thinking of getting from our
        Madame Wang to-day the loan of the garden for the nonce and spreading two
        tables with our mean wine,}
}
{
}

\Verse{224}
{
    \zA{薛姨妈笑道：“果然如此， 算我的孝心虔了。”}
}
{
    \eA{and inviting you, worthy senior, to enjoy the snow; but as I saw that you
        were having a rest, and I heard, at an early hour,}
}
{
}

\Verse{225}
{
    \zA{凤姐儿笑道：“姨妈怎么忘了!如今现秤五十两银子来，交给 我收着，一下雪我就预备下酒。}
}
{
    \eA{that Pao-yü had said that you were not in a joyful frame of mind, I did not,
        in consequence, presume to come and disturb you to-day. But had I known
        sooner the real state of affairs, I would have felt it my bounden duty to
        have asked you round."}
}
{
}

\Verse{226}
{
    \zA{姨妈也不用操心，也不得忘了。”}
}
{
    \eA{"This is," rejoined dowager lady Chia with a smile, "only the first fall of
        snow in the tenth moon.}
}
{
}

\Verse{227}
{
    \zA{贾母笑道：“既 这么说，姨太太给他五十两银子收着，我和他每人分二十五两，到下雪的日子，我 装心里不爽，混过去了。}
}
{
    \eA{We'll have, after this, plenty of snowy days so there will be ample time to
        put your ladyship to wasteful expense." "Verily in that case," Mrs. Hsüeh
        laughingly added, "my filial intentions may well be looked upon as having
        been accomplished." "Mrs. Hsüeh," interposed lady Feng smiling, "mind you
        don't forget it! But you might as well weigh fifty taels this very moment,}
}
{
}

\Verse{228}
{
    \zA{姨太太更不用操心，我和凤姐倒得实惠呢。”}
}
{
    \eA{and hand them over to me to keep, until the first fall of snow, when I can
        get everything ready for the banquet. In this way, you will neither have
        anything to bother you,}
}
{
}

\Verse{229}
{
    \zA{凤姐将手一 拍，笑道：“妙极!这和我的主意一样。”}
}
{
    \eA{aunt, nor will you have a chance of forgetting." "Well, since that be so,"
        old lady Chia remarked with a laugh, "your ladyship had better give her
        fifty taels,}
}
{
}

\Verse{230}
{
    \zA{众人都笑了。}
}
{
    \eA{and I'll share it with her; each one of us taking twenty-five taels; and on
        any day it might snow,}
}
{
}

\Verse{231}
{
    \zA{贾母笑道：“呸!没脸的， 就顺着竿子爬上来了!你不说：姨太太是客，在咱们家受屈，我们该请姨太太才是， 那里有破费姨太太的理？}
}
{
    \eA{I'll pretend I don't feel in proper trim and let it slip by. You'll have
        thus still less occasion to trouble yourself, and I and lady Feng will reap
        a substantial benefit." Lady Feng clapped her hands. "An excellent idea,"
        she laughed. "This quite falls in with my views." The whole company were
        much amused.}
}
{
}

\Verse{232}
{
    \zA{不这么说呢，还有脸先要五十两银子，真不害臊。”}
}
{
    \eA{"Pshaw!" dowager lady Chia laughingly ejaculated. "You barefaced thing!
        (You're like a snake, which) avails itself of the rod, with which it is
        being beaten, to crawl up (and do harm)!}
}
{
}

\Verse{233}
{
    \zA{凤姐 笑道：“我们老祖宗最是有眼色的，试一试姨妈：要松呢，拿出五十两来，就和我 分；这会子估量着不中用了，翻过来拿我做法子，说出这些大方话来。}
}
{
    \eA{You don't try to convince us that it properly devolves upon us, as Mrs.
        Hsüeh is our guest and receives such poor treatment in our household, to
        invite her; for with what right could we subject her ladyship to any
        reckless outlay? but you have the impudence, of impressing upon our minds to
        insist upon the payment, in advance, of fifty taels! Are you really not
        thoroughly ashamed of yourself?"}
}
{
}

\Verse{234}
{
    \zA{如今我也不 和姨妈要银子了，我竟替姨妈出银子，治了酒，请老太太吃了，我另外再封五十两 银子孝敬老祖宗，算是罚我个包揽闲事，这可好不好？”}
}
{
    \eA{"Oh, worthy senior," lady Feng laughed, "you're most sharp-sighted! You try
        to see whether Mrs. Hsüeh will be soft enough to produce fifty taels for you
        to share with me, but fancying now that it's of no avail, you turn round and
        begin to rate me by coming out with all these grand words! I won't however
        take any money from you, Mrs. Hsüeh. I'll, in fact, contribute some on your
        ladyship's account,}
}
{
}

\Verse{235}
{
    \zA{话未说完，众人都笑倒在 炕上。}
}
{
    \eA{and when I get the banquet ready and invite you, venerable ancestor, to come
        and partake of it, I'll also wrap fifty taels in a piece of paper, and
        dutifully present them to you,}
}
{
}

\Verse{236}
{
    \zA{贾母因又说及宝琴雪下折梅，比画儿上还好；又细问他的年庚八字并家内景 况。}
}
{
    \eA{as a penalty for my officious interference in matters that don't concern me.
        Will this be all right or not?" Before these words were brought to a close,
        the various inmates were so convulsed with hearty laughter that they reeled
        over on the stove-couch. Dowager lady Chia then went on to explain how much
        nicer Pao-ch'in was,}
}
{
}

\Verse{237}
{
    \zA{薛姨妈度其意思，大约是要给他求配。}
}
{
    \eA{plucking plum blossom in the snow, than the very picture itself; and she
        next minutely inquired what the year,}
}
{
}

\Verse{238}
{
    \zA{薛姨妈心中因也遂意，只是已许过梅家 了，因贾母尚未说明，自己也不拟定，遂半吐半露告诉贾母道：“可惜了这孩子没 福，前年他父亲就没了。}
}
{
    \eA{moon, day and hour of her birth were, and how things were getting on in her
        home. Mrs. Hsüeh conjectured that the object she had in mind was, in all
        probability, to seek a partner for her. In the secret recesses of her heart,
        Mrs. Hsüeh on this account fell in also with her views.}
}
{
}

\Verse{239}
{
    \zA{他从小儿见的世面倒多，跟他父亲四山五岳都走遍了。}
}
{
    \eA{(Pao-ch'in) had, however, already been promised in marriage to the Mei
        family. But as dowager lady Chia had made, as yet, no open allusion to her
        intentions,}
}
{
}

\Verse{240}
{
    \zA{他 父亲好乐的，各处因有买卖，带了家眷这一省逛一年，明年又到那一省逛半年，所 以天下十停走了有五六停了，那年在这里，把他许了梅翰林的儿子，偏第二年他父
        亲就辞世了。}
}
{
    \eA{(Mrs. Hsüeh) did not think it nice on her part to come out with any definite
        statement, and she accordingly observed to old lady Chia in a vague sort of
        way: "What a pity it is that this girl should have had so little good
        fortune as to lose her father the year before last. But ever since her youth
        up, she has seen much of the world, for she has been with her parent to
        every place of note.}
}
{
}

\Verse{241}
{
    \zA{如今他母亲又是痰症。”}
}
{
    \eA{Her father was a man fond of pleasure; and as he had business in every
        direction,}
}
{
}

\Verse{242}
{
    \zA{凤姐儿也不等说完，便？}
}
{
    \eA{he took his family along with him. After tarrying in this province for a
        whole year,}
}
{
}

\Verse{243}
{
    \zA{声跺脚的说：“偏 不巧!我正要做个媒呢，又已经许了人家！”}
}
{
    \eA{he would next year again go to that province, and spend half a year roaming
        about it everywhere. Hence it is that he had visited five or six tenths of
        the whole empire.}
}
{
}

\Verse{244}
{
    \zA{贾母笑道：“你要给谁说媒？”}
}
{
    \eA{The other year, when they were here, he engaged her to the son of the Hanlin
        Mei.}
}
{
}

\Verse{245}
{
    \zA{凤姐 儿笑道：“老祖宗别管。}
}
{
    \eA{But, as it happened, her father died the year after, and here is her mother
        too now ailing from a superfluity of phlegm."}
}
{
}

\Verse{246}
{
    \zA{心里看准了，他们两个是一对。}
}
{
    \eA{Lady Feng gave her no time to complete what she meant to say. "Hai!" she
        exclaimed, stamping her foot.}
}
{
}

\Verse{247}
{
    \zA{如今有了人家，说也无益， 不如不说罢了。”}
}
{
    \eA{"What you say isn't opportune! I was about to act as a go-between. But is
        she too already engaged?" "For whom did you mean to act as go-between?"}
}
{
}

\Verse{248}
{
    \zA{贾母也知凤姐儿的意思，听见已有人家，也就不提了。}
}
{
    \eA{old lady Chia smiled. "My dear ancestor," lady Feng remarked, "don't concern
        yourself about it! I had determined in my mind that those two would make a
        suitable match.}
}
{
}

\Verse{249}
{
    \zA{大家又闲 话了一会方散。}
}
{
    \eA{But as she has now long ago been promised to some one, it would be of no
        use, were I even to speak out.}
}
{
}

\Verse{250}
{
    \zA{一宿无话。}
}
{
    \eA{Isn't it better that I should hold my peace,}
}
{
}

\Verse{251}
{
    \zA{次日雪晴。}
}
{
    \eA{and drop the whole thing?"}
}
{
}

\Verse{252}
{
    \zA{饭后，贾母又吩咐惜春：“不管冷暖，你要画去；赶到年下，十分 不能，就罢了。}
}
{
    \eA{Dowager lady Chia herself was cognizant of lady Feng's purpose, so upon
        hearing that she already had a suitor, she at once desisted from making any
        further reference to the subject.}
}
{
}

\Verse{253}
{
    \zA{第一要紧把昨儿琴儿和丫头、梅花，照样一笔别错快快添上。”}
}
{
    \eA{The whole company then continued another chat on irrelevant matters for a
        time, after which, they broke up. Nothing of any interest transpired the
        whole night.}
}
{
}

\Verse{254}
{
    \zA{惜 春听了，虽是为难的事，就应了。}
}
{
    \eA{The next day, the snowy weather had cleared up. After breakfast, her
        grandmother Chia again pressed Hsi Ch'un.}
}
{
}

\Verse{255}
{
    \zA{一时众人都来看他如何画。}
}
{
    \eA{"You should go on," she said, "with your painting, irrespective of cold or
        heat.}
}
{
}

\Verse{256}
{
    \zA{惜春只是出神。}
}
{
    \eA{If you can't absolutely finish it by the end of the year,}
}
{
}

\Verse{257}
{
    \zA{李纨 因笑向众人道：“让他自己想去，咱们且说话儿。}
}
{
    \eA{it won't much matter! The main thing is that you must at once introduce in
        it Ch'in Erh and the maid with the plum blossom, as we saw them yesterday,}
}
{
}

\Verse{258}
{
    \zA{昨儿老太太只叫做灯谜儿，回到 家和绮儿纹儿睡不着，我就编了两个《四书》的。}
}
{
    \eA{in strict accordance with the original and without the least discrepancy of
        so much as a stroke." Hsi Ch'un listened to her and felt it her duty to
        signify her assent, in spite of the task being no easy one for her to
        execute.}
}
{
}

\Verse{259}
{
    \zA{他两个每人也编了两个。”}
}
{
    \eA{After a time, a number of her relatives came, in a body, to watch the
        progress of the painting.}
}
{
}

\Verse{260}
{
    \zA{众人 听了，都笑道：“这倒该做的。}
}
{
    \eA{But they discovered Hsi Ch'un plunged in a reverie. "Let's leave her alone,"
        Li Wan smilingly observed to them all,}
}
{
}

\Verse{261}
{
    \zA{先说了，我们猜猜。”}
}
{
    \eA{"to proceed with her meditations; we can meanwhile have a chat among
        ourselves.}
}
{
}

\Verse{262}
{
    \zA{李纨笑道：“‘观音未有世 家传’，打《四书》一句。”}
}
{
    \eA{Yesterday our worthy senior bade us devise a few lantern-conundrums, so when
        we got home, I and Ch'i Erh and Wen Erh did not turn in (but set to work).}
}
{
}

\Verse{263}
{
    \zA{湘云接着就说道：“‘在止于至善’。”}
}
{
    \eA{I composed a couple on the Four Books; but those two girls also managed to
        put together another pair of them."}
}
{
}

\Verse{264}
{
    \zA{宝钗笑道： “你也想一想‘世家传’三个字的意思再猜。”}
}
{
    \eA{"We should hear what they're like," they laughingly exclaimed in chorus,
        when they heard what they had done. "Tell them to us first, and let's have a
        guess!"}
}
{
}

\Verse{265}
{
    \zA{李纨笑道：“再想。”}
}
{
    \eA{"The goddess of mercy has not been handed down by any ancestors."}
}
{
}

\Verse{266}
{
    \zA{黛玉笑道： “我猜罢。}
}
{
    \eA{Li Ch'i smiled. "This refers to a passage in the Four Books."}
}
{
}

\Verse{267}
{
    \zA{可是‘虽善无征’？”}
}
{
    \eA{"In one's conduct, one must press towards the highest benevolence."}
}
{
}

\Verse{268}
{
    \zA{众人都笑道：“这句是了。”}
}
{
    \eA{Hsiang-yün quickly interposed; taking up the thread of the conversation.}
}
{
}

\Verse{269}
{
    \zA{李纨又道：“‘一 池青草草何名’。”}
}
{
    \eA{"You should ponder over the meaning of the three words implying: 'handed
        down by ancestors'," Pao-ch'ai smiled, "before you venture a guess."}
}
{
}

\Verse{270}
{
    \zA{湘云又忙道：“这一定是‘蒲芦也’，再不是不成？”}
}
{
    \eA{"Think again!" Li Wan urged with a smile. "I've guessed it!" Tai-yü smiled.
        "It's: "'If, notwithstanding all that benevolence, there be no outward
        visible sign…'" "That's the line,"}
}
{
}

\Verse{271}
{
    \zA{李纨笑 道：“这难为你猜。}
}
{
    \eA{one and all unanimously exclaimed with a laugh. "'The whole pond is covered
        with rush.'"}
}
{
}

\Verse{272}
{
    \zA{纹儿的是‘水向石边流出冷’，打一古人名。”}
}
{
    \eA{"Now find the name of the rush?" Li Wan proceeded. "This must certainly be
        the cat-tail rush!" hastily again replied Hsiang-yün.}
}
{
}

\Verse{273}
{
    \zA{探春笑着问道： “可是山涛？”}
}
{
    \eA{"Can this not be right?" "You've succeeded in guessing it," Li Wan smiled.
        "Li Wen's is:}
}
{
}

\Verse{274}
{
    \zA{李纨道：“是。”}
}
{
    \eA{"'Cold runs the stream along the stones;'}
}
{
}

\Verse{275}
{
    \zA{李纨又道：“绮儿是个‘萤’字，打一个字。”}
}
{
    \eA{"bearing on the name of a man of old." "Can it be Shan T'ao?" T'an Ch'un
        smilingly asked. "It is!" answered Li Wan. "Ch'i Erh's is the character
        'Yung' (glow-worm).}
}
{
}

\Verse{276}
{
    \zA{众人猜了半日，宝琴道：“这个意思却深，不知可是花草的‘花’字？”}
}
{
    \eA{It refers to a single word," Li Wan resumed. The party endeavoured for a
        long time to hit upon the solution. "The meaning of this is certainly deep,"
        Pao-ch'in put in. "I wonder whether it's the character,}
}
{
}

\Verse{277}
{
    \zA{李绮笑道： “恰是了。”}
}
{
    \eA{'hua,' (flower) in the combination, 'hua ts'ao, (vegetation)."}
}
{
}

\Verse{278}
{
    \zA{众人道：“萤与花何干？”}
}
{
    \eA{"That's just it!" Li Ch'i smiled. "What has a glow-worm to do with flowers?"}
}
{
}

\Verse{279}
{
    \zA{黛玉笑道：“妙的很，萤可不是草化的？”}
}
{
    \eA{one and all observed. "It's capital!" Tai-yü ventured with a smile. "Isn't a
        glow-worm transformed from plants?"}
}
{
}

\Verse{280}
{
    \zA{众人会意，都笑了，说：“好。”}
}
{
    \eA{The company grasped the sense; and, laughing the while,}
}
{
}

\Verse{281}
{
    \zA{宝钗道：“这些虽好，不合老太太的意。}
}
{
    \eA{they, with one consent, shouted out, "splendid!" "All these are, I admit,
        good," Pao-ch'ai remarked,}
}
{
}

\Verse{282}
{
    \zA{不如做些浅近的物儿，大家雅俗共赏 才好。”}
}
{
    \eA{"but they won't suit our venerable senior's taste. Won't it be better
        therefore to compose a few on some simple objects;}
}
{
}

\Verse{283}
{
    \zA{众人都道：“也要做些浅近的俗物才是。”}
}
{
    \eA{some which all of us, whether polished or unpolished, may be able to enjoy?"
        "Yes,"}
}
{
}

\Verse{284}
{
    \zA{湘云想了一想，笑道：“我编 了一支《点绛唇》，却真是个俗物，你们猜猜。”}
}
{
    \eA{they all replied, "we should also think of some simple ones on ordinary
        objects." "I've devised one on the 'Tien Chiang Ch'un' metre," Hsiang-yün
        pursued, after some reflection.}
}
{
}

\Verse{285}
{
    \zA{说着，便念道： 溪壑分离，红尘游戏，真何趣？}
}
{
    \eA{"But it's really on an ordinary object. So try and guess it." Saying this,
        she forthwith went on to recite:}
}
{
}

\Verse{286}
{
    \zA{名利犹虚，后事终难继。}
}
{
    \eA{The creeks and valleys it leaves; Travelling the world, it performs.}
}
{
}

\Verse{287}
{
    \zA{众人都不解，想了半日，也有猜是和尚的，也有猜是道士的，也有猜是偶戏人的。}
}
{
    \eA{In truth how funny it is! But renown and gain are still vain; Ever hard
        behind it is its fate. A conundrum. None of those present could fathom what
        it could be.}
}
{
}

\Verse{288}
{
    \zA{宝玉笑了半日道：“都不是。}
}
{
    \eA{After protracted thought, some made a guess, by saying it was a bonze.}
}
{
}

\Verse{289}
{
    \zA{我猜着了，必定是耍的猴儿。”}
}
{
    \eA{Others maintained that it was a Taoist priest. Others again divined that it
        was a marionette.}
}
{
}

\Verse{290}
{
    \zA{湘云笑道：“正是这 个了。”}
}
{
    \eA{"All your guesses are wrong," Pao-yü chimed in, after considerable
        reflection.}
}
{
}

\Verse{291}
{
    \zA{众人道：“前头都好，末后一句怎么样解？”}
}
{
    \eA{"I've got it! It must for a certainty be a performing monkey." "That's
        really it!" Hsiang-yün laughed. "The first part is all right,"}
}
{
}

\Verse{292}
{
    \zA{湘云道：“那一个耍的猴儿 不是剁了尾巴去的？”}
}
{
    \eA{the party observed, "but how do you explain the last line?" "What performing
        monkey," Hsiang-yün asked, "has not had its tail cut off?"}
}
{
}

\Verse{293}
{
    \zA{众人听了都笑起来，说：“偏他编个谜儿也是刁钻古怪的。”}
}
{
    \eA{Hearing this, they exploded into a fit of merriment. "Even," they argued,
        "the very riddles she improvises are perverse and strange!"}
}
{
}

\Verse{294}
{
    \zA{李纨道：“昨日姨妈说，琴妹妹见得世面多，走的道路也多，你正该编谜儿。}
}
{
    \eA{"Mrs. Hsüeh mentioned yesterday that you, cousin Ch'in, had seen much of the
        world," Li Wan put in, "and that you had also gone about a good deal.}
}
{
}

\Verse{295}
{
    \zA{况且你的诗又好，为什么不编几个儿我们猜一猜？”}
}
{
    \eA{It's for you therefore to try your hand at a few conundrums. What's more
        your poetry too is good. So why shouldn't you indite a few for us to guess?"}
}
{
}

\Verse{296}
{
    \zA{宝琴听了，点头含笑，自去寻 思。}
}
{
    \eA{Pao-ch'in, at this proposal, nodded her head, and while repressing a smile,
        she went off by herself to give way to thought.}
}
{
}

\Verse{297}
{
    \zA{宝钗也有一个，念道： 镂檀镌梓一层层，岂系良工堆砌成？}
}
{
    \eA{Pao-ch'ai then also gave out this riddle: Carved sandal and cut cedar rise
        layer upon layer. Have they been piled and fashioned by workmen of skill!}
}
{
}

\Verse{298}
{
    \zA{虽是半天风雨过，何曾闻得梵铃声？}
}
{
    \eA{In the mid-heavens it's true, both wind and rain fleet by; But can one hear
        the tingling of the Buddhists' bell?}
}
{
}

\Verse{299}
{
    \zA{众人猜时，宝玉也有一个，念道： 天上人间两渺茫，琅？}
}
{
    \eA{While they were giving their mind to guessing what it could be, Pao-yü too
        recited: Both from the heavens and from the earth, it's indistinct to view.}
}
{
}

\Verse{300}
{
    \zA{节过谨提防。}
}
{
    \eA{What time the 'Lang Ya' feast goes past,}
}
{
}

\Verse{301}
{
    \zA{鸾音鹤信须凝睇，好把唏嘘答上苍。}
}
{
    \eA{then mind you take great care. When the 'luan's' notes you catch and the
        crane's message thou'lt look up:}
}
{
}

\Verse{302}
{
    \zA{黛玉也有了一个，念道： ？？}
}
{
    \eA{It is a splendid thing to turn and breathe towards the vault of heaven, (a
        kite) Tai-yü next added:}
}
{
}

\Verse{303}
{
    \zA{何劳缚紫绳？}
}
{
    \eA{Why need a famous steed be a with bridle e'er restrained?}
}
{
}

\Verse{304}
{
    \zA{驰城逐堑势狰狞。}
}
{
    \eA{Through the city it speeds; the moat it skirts; how fierce it looks.}
}
{
}

\Verse{305}
{
    \zA{主人指示风云动，鳌背三山独立名。}
}
{
    \eA{The master gives the word and wind and clouds begin to move. On the 'fish
        backs' and the 'three isles' it only makes a name,}
}
{
}

\Verse{306}
{
    \zA{探春也有了一个，方欲念时，宝琴走来，笑道：“从小儿所走的地方的古迹不 少，我也来挑了十个地方古迹，做了十首‘怀古诗’。}
}
{
    \eA{(a rotating lantern). T'an Ch'un had also one that she felt disposed to tell
        them, but just as she was about to open her lips, Pao-ch'in walked up to
        them. "The relics of various places I've seen since my youth," she smiled,
        "are not few, so I've now selected ten places of historic interest,}
}
{
}

\Verse{307}
{
    \zA{诗虽粗鄙，却怀往事，又暗 隐俗物十件，姐姐们请猜一猜。”}
}
{
    \eA{on which I've composed ten odes, treating of antiquities. The verses may
        possibly be coarse, but they bear upon things of the past, and secretly
        refer as well to ten commonplace articles.}
}
{
}

\Verse{308}
{
    \zA{众人听了，都说：“这倒巧，何不写出来大家一 看？”}
}
{
    \eA{So, cousins, please try and guess them!" "This is ingenious!" they exclaimed
        in chorus, when they heard the result of her labour. "Why not write them
        out,}
}
{
}

\Verse{309}
{
    \zA{要知端的，且听下回分解。}
}
{
    \eA{and let us have a look at them?" But, reader, peruse the next chapter,}
}
{
}

\Verse{310}
{
    \zA{(本章完)}
}
{
    \eA{if you want to learn what follows.}
}
{
}
